% =====================================================================
%  Preface
%  Gattringer & Lang, "Quantum Chromodynamics on the Lattice" (LNP 788)
%  Reconstructed from extract/preface.txt.
% =====================================================================

\chapter*{Preface}
\addcontentsline{toc}{chapter}{Preface}
\markboth{Preface}{Preface}

Quantum chromodynamics (QCD) is the fundamental quantum field theory
of quarks and gluons. In order to discuss it in a mathematically
well-defined way, the theory has to be regularized. Replacing
space--time by a Euclidean lattice has proven to be an efficient
approach which allows for both theoretical understanding and
computational analysis. Lattice QCD has become a standard tool in
elementary particle physics.

As the title already says: this book is introductory! The text is
intended for newcomers to the field, serving as a starting point. We
simply wanted to have a book which we can put into the hands of an
advanced student for a first reading on lattice QCD. This imaginary
student brings as a prerequisite knowledge of higher quantum
mechanics, some continuum quantum field theory, and basic facts of
elementary particle physics phenomenology.

In view of the wealth of applications in current research the topics
presented here are limited and we had to make some painful choices. We
discuss QCD but omit most other lattice field theory applications like
scalar theories, gauge--Higgs models, or electroweak theory. Although
we try to lead the reader up to present day understanding, we cannot
possibly address all ongoing activities, in particular concerning the
role of QCD in electroweak theory. Subjects like glueballs, topological
excitations, and approaches like chiral perturbation theory are
mentioned only briefly. This allows us to cover the other topics quite
explicitly, including detailed derivations of key equations. The field
is rapidly developing. The proceedings of the annual lattice
conferences provide information on newer directions and up-to-date
results.

As usual, completing the book took longer than originally planned and
we thank our editor Claus Ascheron for his patience. We are very
grateful to many of our colleagues, who offered to read one or the
other piece. In particular we want to thank Vladimir Braun,
Dirk Br\"ommel, Tommy Burch, Stefano Capitani, Tom DeGrand,
Stephan D\"urr, Georg Engel, Christian Hagen, Leonid Glozman,
Meinulf G\"ockeler, Peter Hasenfratz, Jochen Heitger, Verena Hermann,
Edwin Laermann, Markus Limmer, Pushan Majumdar, Daniel Mohler,
Wolfgang Ortner, Bernd-Jochen Schaefer, Stefan Schaefer, Andreas
Sch\"afer, Erhard Seiler, Stefan Sint, Stefan Solbrig, and Pierre van
Baal.

It would be surprising if there were not mistakes in this text. We
therefore set up a web companion to this book: \url{http://physik.uni-graz.at/qcdlatt/}
On that page we document errata and provide further links and
information.

\begin{flushright}
Graz, March 2009\\[2mm]
Christof Gattringer\\[1mm]
Christian B. Lang
\end{flushright}
